% ============================================================================
% CAPÍTULO 5: ESTIMAÇÃO DE PARÂMETROS
% Bioinformática para Biologia de Sistemas
% ============================================================================

% ============================================================================
% TEMPLATE BASE PARA APRESENTAÇÕES EM BEAMER
% Bioinformática para Biologia de Sistemas
% ============================================================================

\documentclass[12pt, aspectratio=169]{beamer}

% ============================================================================
% PACOTES ESSENCIAIS
% ============================================================================
\usepackage[utf8]{inputenc}
\usepackage[T1]{fontenc}
\usepackage[brazilian]{babel}
\usepackage{amsmath, amssymb, amsthm}
\usepackage{graphicx}
\usepackage{xcolor}
\usepackage{tikz}
\usetikzlibrary{arrows.meta, shapes, positioning, calc, decorations.pathreplacing, decorations.shapes, decorations.pathmorphing}
\usepackage{listings}
\usepackage{booktabs}
\usepackage{multirow}
\usepackage{hyperref}
\usepackage{chemfig}
\usepackage{chemarr}

% ============================================================================
% CONFIGURAÇÃO DO TEMA
% ============================================================================
\usetheme{Madrid}
\usecolortheme{default}

% Cores personalizadas
\definecolor{azulescuro}{RGB}{0, 51, 102}
\definecolor{azulclaro}{RGB}{51, 153, 255}
\definecolor{verdeacido}{RGB}{102, 204, 0}
\definecolor{laranjacel}{RGB}{255, 153, 51}
\definecolor{roxodna}{RGB}{153, 51, 255}

\setbeamercolor{structure}{fg=azulescuro}
\setbeamercolor{palette primary}{bg=azulescuro,fg=white}
\setbeamercolor{palette secondary}{bg=azulclaro,fg=white}
\setbeamercolor{palette tertiary}{bg=azulescuro,fg=white}
\setbeamercolor{palette quaternary}{bg=azulclaro,fg=white}
\setbeamercolor{block title}{bg=azulescuro,fg=white}
\setbeamercolor{block body}{bg=azulclaro!10,fg=black}
\setbeamercolor{block title alerted}{bg=red!80,fg=white}
\setbeamercolor{block body alerted}{bg=red!10,fg=black}
\setbeamercolor{block title example}{bg=verdeacido!80,fg=white}
\setbeamercolor{block body example}{bg=verdeacido!10,fg=black}

% ============================================================================
% CONFIGURAÇÃO DE CÓDIGO
% ============================================================================
\lstset{
    language=Python,
    basicstyle=\ttfamily\small,
    keywordstyle=\color{azulescuro}\bfseries,
    commentstyle=\color{verdeacido},
    stringstyle=\color{laranjacel},
    numbers=left,
    numberstyle=\tiny\color{gray},
    stepnumber=1,
    numbersep=5pt,
    backgroundcolor=\color{white},
    showspaces=false,
    showstringspaces=false,
    showtabs=false,
    frame=single,
    rulecolor=\color{black},
    tabsize=2,
    captionpos=b,
    breaklines=true,
    breakatwhitespace=false,
    escapeinside={\%*}{*)},
    xleftmargin=10pt,
    xrightmargin=10pt
}

% Definir estilo para R
\lstdefinestyle{Rstyle}{
    language=R,
    basicstyle=\ttfamily\small,
    keywordstyle=\color{azulescuro}\bfseries,
    commentstyle=\color{verdeacido},
    stringstyle=\color{laranjacel}
}

% Definir estilo para MATLAB
\lstdefinestyle{MATLABstyle}{
    language=Matlab,
    basicstyle=\ttfamily\small,
    keywordstyle=\color{azulescuro}\bfseries,
    commentstyle=\color{verdeacido},
    stringstyle=\color{laranjacel}
}

% ============================================================================
% COMANDOS PERSONALIZADOS
% ============================================================================

% Comando para equações destacadas
\newcommand{\destaque}[1]{%
    \begin{alertblock}{}
        \begin{center}
            #1
        \end{center}
    \end{alertblock}
}

% Comando para definições
\newcommand{\definicao}[2]{%
    \begin{block}{Definição: #1}
        #2
    \end{block}
}

% Comando para exemplos
\newcommand{\exemplo}[2]{%
    \begin{exampleblock}{Exemplo: #1}
        #2
    \end{exampleblock}
}

% Comando para observações importantes
\newcommand{\importante}[1]{%
    \begin{alertblock}{Importante!}
        #1
    \end{alertblock}
}

% Comandos matemáticos úteis
\newcommand{\R}{\mathbb{R}}
\newcommand{\N}{\mathbb{N}}
\newcommand{\Z}{\mathbb{Z}}
\newcommand{\dif}{\mathrm{d}}
\newcommand{\parcial}[2]{\frac{\partial #1}{\partial #2}}
\newcommand{\derivada}[2]{\frac{\dif #1}{\dif #2}}
\newcommand{\vect}[1]{\boldsymbol{#1}}

% Comandos biológicos
\newcommand{\gene}[1]{\textit{#1}}
\newcommand{\proteina}[1]{\textsc{#1}}
\newcommand{\especie}[1]{\textit{#1}}

% ============================================================================
% CONFIGURAÇÃO DE RODAPÉ
% ============================================================================
\setbeamertemplate{footline}{
  \leavevmode%
  \hbox{%
  \begin{beamercolorbox}[wd=.33\paperwidth,ht=2.25ex,dp=1ex,center]{author in head/foot}%
    \usebeamerfont{author in head/foot}\insertshortauthor
  \end{beamercolorbox}%
  \begin{beamercolorbox}[wd=.34\paperwidth,ht=2.25ex,dp=1ex,center]{title in head/foot}%
    \usebeamerfont{title in head/foot}\insertshorttitle
  \end{beamercolorbox}%
  \begin{beamercolorbox}[wd=.33\paperwidth,ht=2.25ex,dp=1ex,right]{date in head/foot}%
    \usebeamerfont{date in head/foot}\insertshortdate{}\hspace*{2em}
    \insertframenumber{} / \inserttotalframenumber\hspace*{2ex}
  \end{beamercolorbox}}%
  \vskip0pt%
}

% ============================================================================
% CONFIGURAÇÃO DE NAVEGAÇÃO
% ============================================================================
\setbeamertemplate{navigation symbols}{}

% ============================================================================
% CONFIGURAÇÃO DE ITEMIZE/ENUMERATE
% ============================================================================
\setbeamertemplate{itemize items}[circle]
\setbeamertemplate{enumerate items}[default]

% ============================================================================
% FIM DO TEMPLATE
% ============================================================================


% ============================================================================
% INFORMAÇÕES DO DOCUMENTO
% ============================================================================
\title[Cap. 5: Estimação de Parâmetros]{Capítulo 5: Estimação de Parâmetros}
\subtitle{Bioinformática para Biologia de Sistemas}
\author{Ney Lemke}
\institute{DBF-Unesp}
\date{\today}

\begin{document}

% ============================================================================
% SLIDE DE TÍTULO
% ============================================================================
\begin{frame}
\titlepage
\end{frame}

% ============================================================================
% SUMÁRIO
% ============================================================================
\begin{frame}{Sumário}
\tableofcontents
\end{frame}

% ============================================================================
% SEÇÃO 1: INTRODUÇÃO
% ============================================================================
\section{Introdução}

\begin{frame}{5.1 Por que Estimação de Parâmetros?}
\begin{block}{Objetivos de Aprendizagem}
\begin{itemize}
    \item Compreender o problema inverso em biologia de sistemas
    \item Dominar métodos de otimização local e global
    \item Aplicar técnicas de estimação a modelos biológicos
    \item Avaliar identificabilidade e incerteza de parâmetros
    \item Validar modelos com critérios estatísticos
\end{itemize}
\end{block}

\vspace{0.3cm}

\definicao{Estimação de Parâmetros}{
Processo de determinar valores numéricos de parâmetros de um modelo matemático a partir de dados experimentais, ajustando o modelo para reproduzir o comportamento observado do sistema.
}
\end{frame}

\begin{frame}{O Problema Inverso}
\begin{columns}
\column{0.48\textwidth}
\textbf{Problema Direto:}
\begin{itemize}\small
    \item Parâmetros conhecidos
    \item Condições iniciais dadas
    \item Simular $\rightarrow$ Previsões
    \item Determinístico
\end{itemize}

\vspace{0.08cm}

\begin{exampleblock}{Exemplo}\small
Dados $k_1=0.5$, $k_2=0.1$\\
Simular: $\frac{dx}{dt} = k_1 - k_2 x$\\
$\Rightarrow$ obter $x(t)$
\end{exampleblock}

\column{0.48\textwidth}
\textbf{Problema Inverso:}
\begin{itemize}\small
    \item Dados experimentais
    \item Parâmetros desconhecidos
    \item Estimar $\leftarrow$ Ajuste
    \item Mal-posto (ill-posed)
\end{itemize}

\vspace{0.08cm}

\begin{alertblock}{Exemplo}\small
Dados $x(t)$ experimentais\\
Modelo: $\frac{dx}{dt} = k_1 - k_2 x$\\
$\Rightarrow$ encontrar $k_1$, $k_2$
\end{alertblock}
\end{columns}

\vspace{0.1cm}

\importante{
O problema inverso é geralmente mais difícil que o direto!
}
\end{frame}

\begin{frame}{Desafios na Estimação de Parâmetros}
\begin{columns}
\column{0.55\textwidth}
\begin{block}{Principais Obstáculos}
\begin{enumerate}\small
    \item \textbf{Identificabilidade:} múltiplas soluções
    \item \textbf{Ruído:} medições imprecisas
    \item \textbf{Dados escassos:} poucas observações
    \item \textbf{Mínimos locais:} otimização não-convexa
    \item \textbf{Correlação:} parâmetros colineares
    \item \textbf{Custo computacional:} simulações
    \item \textbf{Incerteza:} quantificação de confiança
\end{enumerate}
\end{block}

\column{0.41\textwidth}
\begin{center}
\begin{tikzpicture}[scale=0.6]
    % Landscape with multiple minima
    \draw[->] (0,0) -- (5,0) node[right] {\tiny $\theta$};
    \draw[->] (0,0) -- (0,2.5) node[above] {\tiny $J(\theta)$};

    % Objective function
    \draw[thick, azulescuro, domain=0.4:4.6, samples=100, smooth]
        plot (\x, {0.8 + 0.4*sin(4*\x r) + 0.3*(\x-2.5)^2/3});

    % Local minima
    \fill[red] (0.9,0.5) circle (1.5pt) node[above, font=\tiny] {Local};
    \fill[red] (3.9,1.1) circle (1.5pt) node[above, font=\tiny] {Local};
    \fill[verdeacido] (2.3,0.3) circle (1.5pt) node[below, font=\tiny] {Global};
\end{tikzpicture}
\end{center}
\end{columns}
\end{frame}

\begin{frame}{Visão Geral dos Métodos}
\begin{columns}
\column{0.50\textwidth}
\begin{center}
\begin{tikzpicture}[scale=0.8, every node/.style={transform shape}, node distance=2.0cm, auto,
    box/.style={rectangle, draw, text width=2.4cm, text centered, rounded corners, minimum height=0.8cm, font=\scriptsize}]

    \node[box, fill=azulclaro!30] (methods) {Métodos de\\Estimação};
    \node[box, fill=laranjacel!30, above left=1.2cm and 0.6cm of methods] (grid) {Busca em\\Grade};
    \node[box, fill=verdeacido!30, above right=1.2cm and 0.6cm of methods] (local) {Otimização\\Local};
    \node[box, fill=roxodna!30, below=1.0cm of methods] (global) {Otimização\\Global};

    \draw[->, thick] (methods) -- (grid);
    \draw[->, thick] (methods) -- (local);
    \draw[->, thick] (methods) -- (global);

    \node[above, font=\tiny] at ($(methods)!0.5!(grid)$) {Simples};
    \node[above, font=\tiny] at ($(methods)!0.5!(local)$) {Rápido};
    \node[left, font=\tiny] at ($(methods)!0.5!(global)$) {Robusto};
\end{tikzpicture}
\end{center}

\column{0.48\textwidth}
\begin{exampleblock}{Escolha do Método}
\begin{itemize}\small
    \item \textbf{Poucos parâmetros + boa estimativa:} locais (rápido)
    \item \textbf{Muitos parâmetros + incerteza:} globais (robusto)
    \item \textbf{Exploração inicial:} busca em grade (simples)
\end{itemize}
\end{exampleblock}
\end{columns}
\end{frame}

% ============================================================================
% SEÇÃO 2: FORMULAÇÃO DO PROBLEMA
% ============================================================================
\section{Formulação do Problema}

\begin{frame}{Função Objetivo: Conceito}
\definicao{Função Objetivo (Cost Function)}{
Medida quantitativa da discrepância entre dados experimentais e previsões do modelo. O objetivo na estimação é encontrar parâmetros que minimizam esta função.
}

\begin{block}{Componentes Principais}
\begin{itemize}
    \item $y_i$: dados experimentais observados no ponto $i$
    \item $\hat{y}_i(\vect{\theta})$: previsão ou simulação gerada pelo modelo
    \item $\vect{\theta}$: vetor de parâmetros a serem ajustados/estimados
    \item $n$: número total de observações experimentais
\end{itemize}
\end{block}
\end{frame}

\begin{frame}{Função Objetivo: Formulação}
A função objetivo quantifica o erro total entre os dados e as previsões:

\destaque{
\begin{equation*}
J(\vect{\theta}) = \text{medida de discrepância entre } \{y_i\} \text{ e } \{\hat{y}_i(\vect{\theta})\}
\end{equation*}
}

\vspace{0.4cm}

\textbf{Problema de Otimização:}
A estimação é formulada matematicamente como a busca pelo argumento minimizador de $J(\vect{\theta})$:
\begin{equation*}
\vect{\theta}^* = \arg\min_{\vect{\theta}} J(\vect{\theta})
\end{equation*}
\end{frame}

\begin{frame}{Mínimos Quadrados Ordinários (OLS)}
\begin{block}{Método Mais Comum}
Minimização da soma dos quadrados dos desvios entre o modelo e os dados:
\end{block}

\destaque{
\begin{equation*}
J(\vect{\theta}) = \text{SSE}(\vect{\theta}) = \sum_{i=1}^{n} [y_i - \hat{y}_i(\vect{\theta})]^2
\end{equation*}
}

\vspace{0.2cm}
\begin{exampleblock}{Resíduo}
O resíduo $r_i$ representa o erro de previsão individual no ponto $i$:
\begin{equation*}
r_i = y_i - \hat{y}_i(\vect{\theta}) \quad \Rightarrow \quad \text{SSE}(\vect{\theta}) = \sum_{i=1}^{n} r_i^2
\end{equation*}
\end{exampleblock}
\end{frame}

\begin{frame}{Mínimos Quadrados Ordinários (OLS): Propriedades}
A formulação OLS clássica possui características importantes:

\begin{columns}
\column{0.48\textwidth}
\textbf{Vantagens:}
\begin{itemize}\small
    \item Simples e intuitivo
    \item Bem estabelecido na literatura
    \item Solução analítica para modelos lineares
    \item Interpretação estatística direta
\end{itemize}

\column{0.48\textwidth}
\textbf{Suposições Estatísticas:}
\begin{itemize}\small
    \item Erros gaussianos (ruído branco)
    \item Variância constante (homocedasticidade)
    \item Erros independentes
    \item Outliers são muito problemáticos
\end{itemize}
\end{columns}
\end{frame}

\begin{frame}{Mínimos Quadrados Ponderados (WLS)}
\begin{block}{Quando usar?}
Quando as observações têm diferentes níveis de incerteza ou confiabilidade.
\end{block}

\destaque{
\begin{equation*}
J(\vect{\theta}) = \sum_{i=1}^{n} w_i [y_i - \hat{y}_i(\vect{\theta})]^2
\end{equation*}
}

\begin{columns}
\column{0.52\textwidth}
\textbf{Pesos $w_i$:}
\begin{itemize}\small
    \item $w_i = 1/\sigma_i^2$: inverso da variância (comum)
    \item $w_i = 1/y_i$: proporcional ao valor (dados log-scale)
    \item $w_i$ customizado: baseado no erro experimental
\end{itemize}

\column{0.44\textwidth}
\exemplo{Cinética Enzimática}{
\small
Concentrações baixas têm maior erro relativo.\\
$\Rightarrow$ usar pesos $w_i = 1/\hat{y}_i$ ou $w_i = 1/\hat{y}_i^2$.
}
\end{columns}
\end{frame}

\begin{frame}{Maximum Likelihood Estimation (MLE): Conceito}
\begin{block}{Abordagem Estatística}
Encontrar parâmetros que maximizam a probabilidade (likelihood) de observar os dados reais obtidos.
\end{block}

\textbf{Função de Verossimilhança:}
\begin{equation*}
\mathcal{L}(\vect{\theta}) = P(\text{dados} | \vect{\theta})
\end{equation*}

\textbf{Log-verossimilhança (computacionalmente preferível):}
\begin{equation*}
\ln \mathcal{L}(\vect{\theta}) = \sum_{i=1}^{n} \ln P(y_i | \vect{\theta})
\end{equation*}
\end{frame}

\begin{frame}{Maximum Likelihood Estimation (MLE): Caso Gaussiano}
\begin{exampleblock}{Caso Gaussiano}
Assumindo que os erros de medição $\epsilon_i$ são independentes e normalmente distribuídos, $\epsilon_i \sim \mathcal{N}(0, \sigma^2)$:
\begin{equation*}
\ln \mathcal{L}(\vect{\theta}) = -\frac{n}{2}\ln(2\pi\sigma^2) - \frac{1}{2\sigma^2}\sum_{i=1}^{n}[y_i - \hat{y}_i(\vect{\theta})]^2
\end{equation*}
\end{exampleblock}

\vspace{0.3cm}
\importante{
Sob a hipótese de ruído gaussiano, maximizar a log-verossimilhança $\ln \mathcal{L}(\vect{\theta})$ é matematicamente equivalente a minimizar a soma dos quadrados dos resíduos (SSE) do método de OLS!
}
\end{frame}

\begin{frame}{Abordagem Bayesiana: Teorema e Termos}
\begin{block}{Incorpora Conhecimento Prévio}
Usa o Teorema de Bayes para atualizar a nossa crença sobre as distribuições dos parâmetros à medida que novos dados são observados.
\end{block}

\destaque{
\begin{equation*}
P(\vect{\theta} | \text{dados}) = \frac{P(\text{dados} | \vect{\theta}) \cdot P(\vect{\theta})}{P(\text{dados})}
\end{equation*}
}

\begin{itemize}
    \item $P(\vect{\theta} | \text{dados})$: \textbf{posterior} (distribuição pós-dados)
    \item $P(\text{dados} | \vect{\theta})$: \textbf{likelihood} (verossimilhança dos dados)
    \item $P(\vect{\theta})$: \textbf{prior} (conhecimento prévio sobre os parâmetros)
    \item $P(\text{dados})$: evidência (constante de normalização)
\end{itemize}
\end{frame}

\begin{frame}{Abordagem Bayesiana: Vantagens e Desafios}
\begin{exampleblock}{Vantagens}
\begin{itemize}
    \item Permite incorporar dados históricos ou da literatura nas distribuições priors
    \item Quantifica incerteza de forma natural (probabilidade dos parâmetros)
    \item Previne overfitting agindo como uma regularização natural
\end{itemize}
\end{exampleblock}

\vspace{0.3cm}

\begin{alertblock}{Desafio}
Requer a escolha de distribuições priors adequadas e demanda alto custo computacional com algoritmos de amostragem (como MCMC).
\end{alertblock}
\end{frame}

% ============================================================================
% SEÇÃO 3: BUSCA EM GRADE
% ============================================================================
\section{Busca em Grade (Grid Search)}

\begin{frame}{Conceito de Busca em Grade}
\definicao{Grid Search}{
Método de otimização que avalia sistematicamente a função objetivo em uma grade regular de valores de parâmetros.
}

\begin{columns}
\column{0.48\textwidth}
\textbf{Algoritmo:}
\begin{enumerate}\small
    \item Definir limites para os parâmetros
    \item Criar grade regular com $n_1 \times n_2 \times \ldots$ pontos
    \item Avaliar a função objetivo $J(\vect{\theta})$ em cada nó da grade
    \item Selecionar o ponto com o menor valor de $J$
\end{enumerate}

\column{0.48\textwidth}
\begin{center}
\begin{tikzpicture}[scale=0.7]
    % 2D grid
    \draw[step=0.6cm, gray, very thin] (0,0) grid (3,3);
    \draw[->] (0,0) -- (3.5,0) node[right] {\tiny $\theta_1$};
    \draw[->] (0,0) -- (0,3.5) node[above] {\tiny $\theta_2$};

    % Grid points with values
    \foreach \x in {0.3, 0.9, 1.5, 2.1, 2.7} {
        \foreach \y in {0.3, 0.9, 1.5, 2.1, 2.7} {
            \pgfmathsetmacro{\val}{100*exp(-(\x-1.5)^2/2 - (\y-1.8)^2/2)}
            \pgfmathparse{\val>50 ? 1 : 0}
            \ifnum\pgfmathresult=1
                \fill[verdeacido] (\x,\y) circle (2pt);
            \else
                \pgfmathparse{\val>20 ? 1 : 0}
                \ifnum\pgfmathresult=1
                    \fill[laranjacel] (\x,\y) circle (2pt);
                \else
                    \fill[red] (\x,\y) circle (2pt);
                \fi
            \fi
        }
    }

    % Best point
    \fill[blue] (1.5,1.8) circle (3pt) node[right, font=\tiny] {Mínimo};
    \node[below, font=\tiny] at (1.5,-0.2) {Grade 5x5 = 25 avaliações};
\end{tikzpicture}
\end{center}
\end{columns}
\end{frame}

\begin{frame}{Vantagens e Desvantagens}
\begin{columns}
\column{0.48\textwidth}
\begin{block}{Vantagens}\small
\begin{itemize}
    \item Simples de implementar
    \item Facilmente paralelizável
    \item Garantia de encontrar o mínimo global
    \item Sem derivadas necessárias
\end{itemize}
\end{block}

\column{0.48\textwidth}
\begin{alertblock}{Desvantagens}\small
\begin{itemize}
    \item Custo exponencial: $O(n^p)$
    \item Inviável para muitos parâmetros
    \item Resolução da grade limitada
    \item Maldição da dimensionalidade
\end{itemize}
\end{alertblock}
\end{columns}
\end{frame}

\begin{frame}{Grid Search: Custo Computacional}
O custo computacional cresce exponencialmente com o número de parâmetros estimáveis.

\importante{
\textbf{Custo Computacional:} assumindo 10 valores testados por parâmetro:
\begin{itemize}
    \item 2 parâmetros: $10^2 = 100$ avaliações de $J$
    \item 5 parâmetros: $10^5 = 100.000$ avaliações de $J$
    \item 10 parâmetros: $10^{10}$ avaliações (impraticável em tempo viável!)
\end{itemize}
}

\vspace{0.3cm}
\textbf{Conclusão:} A busca em grade é útil para exploração inicial de poucos parâmetros (1 ou 2), mas inadequada para problemas de alta dimensão.
\end{frame}

\begin{frame}[fragile]{Implementação: Grid Search 1D}
\begin{lstlisting}[language=Python, caption=Busca em grade unidimensional, basicstyle=\ttfamily\scriptsize]
import numpy as np

def objective_function(theta):
    return (theta - 3)**2 + 5

# Definir grade e avaliar
theta_values = np.linspace(0, 6, 50)
J_values = [objective_function(t) for t in theta_values]

# Encontrar minimo
idx_min = np.argmin(J_values)
theta_best = theta_values[idx_min]
J_best = J_values[idx_min]

print(f"Parametro otimo: {theta_best:.3f}")
print(f"Valor da funcao: {J_best:.3f}")
\end{lstlisting}
\end{frame}

\begin{frame}[fragile]{Implementação: Grid Search 2D}
\begin{lstlisting}[language=Python, caption=Busca em grade bidimensional, basicstyle=\ttfamily\scriptsize]
def sse_model(params, t, y):
    a, b = params
    return np.sum((y - a * np.exp(-b * t))**2)

# Dados experimentais e grade
t = np.array([0, 1, 2, 3, 4, 5])
y = np.array([10, 6.1, 3.7, 2.2, 1.4, 0.8])
a_range = np.linspace(5, 15, 20)
b_range = np.linspace(0.1, 1.0, 20)
SSE_grid = np.array([[sse_model([a, b], t, y) for b in b_range] for a in a_range])

# Encontrar minimo
idx = np.unravel_index(SSE_grid.argmin(), SSE_grid.shape)
a_best, b_best = a_range[idx[0]], b_range[idx[1]]
print(f"Parametros otimos: a={a_best:.3f}, b={b_best:.3f}")
\end{lstlisting}
\end{frame}

\begin{frame}[fragile]{Visualização: Contour Plot}
\begin{lstlisting}[language=Python, caption=Mapa de contorno da funcao objetivo, basicstyle=\ttfamily\scriptsize]
# Criar contour plot
A, B = np.meshgrid(a_range, b_range)

plt.figure(figsize=(8, 5))
contour = plt.contourf(A, B, SSE_grid.T, levels=20, cmap='viridis')
plt.colorbar(contour, label='SSE')

# Marcar minimo e contornos
plt.plot(a_best, b_best, 'r*', markersize=15,
         label=f'Minimo: ({a_best:.2f}, {b_best:.2f})')
plt.contour(A, B, SSE_grid.T, levels=10,
            colors='white', linewidths=0.5, alpha=0.5)

plt.xlabel('Parametro a')
plt.ylabel('Parametro b')
plt.title('Landscape da Funcao Objetivo')
plt.legend()
plt.grid(True, alpha=0.3)
\end{lstlisting}
\end{frame}

% ============================================================================
% SEÇÃO 4: MÉTODOS DE GRADIENTE
% ============================================================================
\section{Métodos de Gradiente}

\begin{frame}{Gradient Descent (Descida do Gradiente)}
\definicao{Gradient Descent}{
Método iterativo que se move na direção oposta ao gradiente para encontrar um mínimo local da função objetivo.
}

\destaque{
\begin{equation*}
\vect{\theta}^{(k+1)} = \vect{\theta}^{(k)} - \alpha \nabla J(\vect{\theta}^{(k)})
\end{equation*}
}

\begin{columns}
\column{0.54\textwidth}
\textbf{Componentes:}
\begin{itemize}\small
    \item $\vect{\theta}^{(k)}$: parâmetros na iteração $k$
    \item $\alpha > 0$: taxa de aprendizado (passo)
    \item $\nabla J = \left[\frac{\partial J}{\partial \theta_1}, \frac{\partial J}{\partial \theta_2}, \ldots\right]^T$: vetor gradiente
\end{itemize}

\column{0.42\textwidth}
\begin{center}
\begin{tikzpicture}[scale=0.65]
    \draw[->] (0,0) -- (4.5,0) node[right] {\tiny $\theta$};
    \draw[->] (0,0) -- (0,2.5) node[above] {\tiny $J$};

    % Objective function
    \draw[thick, azulescuro, domain=0.3:4.2, samples=100, smooth]
        plot (\x, {0.3*(\x-2.2)^2 + 0.4});

    % Gradient descent steps
    \foreach \x/\y in {3.6/1.8, 3.1/1.3, 2.7/0.9, 2.4/0.6, 2.3/0.45} {
        \fill[red] (\x,\y) circle (1.5pt);
        \draw[->, thick, verdeacido] (\x,\y) -- +(-.3,-.2);
    }

    \node[above, font=\tiny] at (3.6,1.9) {Início};
    \node[below, font=\tiny] at (2.2,0.3) {Mínimo};
\end{tikzpicture}
\end{center}
\end{columns}
\end{frame}

\begin{frame}{Escolha da Taxa de Aprendizado: Impacto}
\begin{columns}
\column{0.32\textwidth}
\begin{center}
\begin{tikzpicture}[scale=0.6]
    \draw[thick, azulescuro, domain=0:3, samples=50]
        plot (\x, {0.5*(\x-1.5)^2 + 0.3});
    \draw[->, red] (2.5,1) -- (1.8,0.5);
    \draw[->, red] (1.8,0.5) -- (1.55,0.35);
    \node[below, font=\tiny] at (1.5,-0.3) {$\alpha$ pequeno};
    \node[below, font=\tiny] at (1.5,-0.6) {Lento};
\end{tikzpicture}
\end{center}

\column{0.32\textwidth}
\begin{center}
\begin{tikzpicture}[scale=0.6]
    \draw[thick, azulescuro, domain=0:3, samples=50]
        plot (\x, {0.5*(\x-1.5)^2 + 0.3});
    \draw[->, verdeacido] (2.5,1) -- (1.5,0.3);
    \node[below, font=\tiny] at (1.5,-0.3) {$\alpha$ ideal};
    \node[below, font=\tiny] at (1.5,-0.6) {Rápido};
\end{tikzpicture}
\end{center}

\column{0.32\textwidth}
\begin{center}
\begin{tikzpicture}[scale=0.6]
    \draw[thick, azulescuro, domain=0:3, samples=50]
        plot (\x, {0.5*(\x-1.5)^2 + 0.3});
    \draw[->, red] (2.5,1) -- (0.3,1.2);
    \draw[->, red] (0.3,1.2) -- (2.7,1.4);
    \node[below, font=\tiny] at (1.5,-0.3) {$\alpha$ grande};
    \node[below, font=\tiny] at (1.5,-0.6) {Oscila};
\end{tikzpicture}
\end{center}
\end{columns}
\end{frame}

\begin{frame}{Escolha da Taxa de Aprendizado: Estratégias}
\importante{
A escolha de $\alpha$ é crucial para convergência!
}

\begin{block}{Estratégias}
\begin{itemize}
    \item \textbf{Fixo:} $\alpha$ constante (simples, pode não convergir)
    \item \textbf{Decrescente:} $\alpha_k = \alpha_0 / (1+k)$ (garante convergência)
    \item \textbf{Line search:} otimizar $\alpha$ em cada iteração (mais robusto)
    \item \textbf{Adaptativo:} ajustar baseado em progresso (Adam, RMSprop)
\end{itemize}
\end{block}
\end{frame}

\begin{frame}{Critérios de Convergência}
\begin{block}{Quando Parar?}
Teste múltiplos critérios simultaneamente:
\end{block}

\begin{enumerate}
    \item \textbf{Gradiente pequeno:} $\|\nabla J(\vect{\theta}^{(k)})\| < \epsilon_g$
    \begin{itemize}
        \item Indica ponto estacionário
    \end{itemize}

    \item \textbf{Mudança em parâmetros:} $\|\vect{\theta}^{(k+1)} - \vect{\theta}^{(k)}\| < \epsilon_\theta$
    \begin{itemize}
        \item Parâmetros não mudam mais
    \end{itemize}

    \item \textbf{Mudança na função:} $|J^{(k+1)} - J^{(k)}| < \epsilon_J$
    \begin{itemize}
        \item Função objetivo estabilizou
    \end{itemize}

    \item \textbf{Número máximo de iterações:} $k > k_{\max}$
    \begin{itemize}
        \item Prevenir loops infinitos
    \end{itemize}
\end{enumerate}

\begin{exampleblock}{Valores Típicos}
$\epsilon_g = 10^{-6}$, $\epsilon_\theta = 10^{-6}$, $\epsilon_J = 10^{-8}$, $k_{\max} = 1000$
\end{exampleblock}
\end{frame}

\begin{frame}[fragile]{Implementação: Gradient Descent}
\begin{lstlisting}[language=Python, caption=Descida do gradiente principal, basicstyle=\ttfamily\scriptsize]
def gradient_descent(f, theta0, alpha=0.01, max_iter=1000, tol=1e-6):
    theta = np.array(theta0, dtype=float)
    history = [theta.copy()]

    for k in range(max_iter):
        # Calcular gradiente numericamente
        grad = numerical_gradient(f, theta)

        # Atualizar parametros e checar convergencia
        theta_new = theta - alpha * grad
        if np.linalg.norm(theta_new - theta) < tol:
            print(f"Convergiu em {k} iteracoes")
            break

        theta = theta_new
        history.append(theta.copy())

    return theta, np.array(history)
\end{lstlisting}
\end{frame}

\begin{frame}[fragile]{Implementação: Gradiente Numérico}
\begin{lstlisting}[language=Python, caption=Calculo numerico do gradiente, basicstyle=\ttfamily\scriptsize]
def numerical_gradient(f, theta, h=1e-5):
    """Gradiente por diferencas finitas"""
    grad = np.zeros_like(theta)
    for i in range(len(theta)):
        theta_plus = theta.copy()
        theta_plus[i] += h
        grad[i] = (f(theta_plus) - f(theta)) / h
    return grad
\end{lstlisting}
\end{frame}

\begin{frame}{Método de Newton-Raphson}
\begin{block}{Ideia}
Usar informação de segunda ordem (curvatura) para convergência mais rápida
\end{block}

\destaque{
\begin{equation*}
\vect{\theta}^{(k+1)} = \vect{\theta}^{(k)} - \mathbf{H}^{-1}(\vect{\theta}^{(k)}) \nabla J(\vect{\theta}^{(k)})
\end{equation*}
}

\begin{itemize}\small
    \item $\mathbf{H}$: matriz Hessiana (derivadas segundas)
    \item $H_{ij} = \frac{\partial^2 J}{\partial \theta_i \partial \theta_j}$
\end{itemize}

\begin{columns}
\column{0.48\textwidth}
\textbf{Vantagens:}
\begin{itemize}\small
    \item Convergência quadrática
    \item Rápido perto do mínimo
    \item Não precisa de step size
\end{itemize}

\column{0.48\textwidth}
\textbf{Desvantagens:}
\begin{itemize}\small
    \item Calcular $\mathbf{H}$ é caro
    \item Inverter $\mathbf{H}$ também
    \item Não-convergência longe do mínimo
\end{itemize}
\end{columns}
\end{frame}

\begin{frame}{Método de Gauss-Newton}
\begin{block}{Específico para Mínimos Quadrados}\small
Aproxima a Hessiana usando apenas o Jacobiano das previsões do modelo.
\end{block}

Para $J(\vect{\theta}) = \sum r_i^2(\vect{\theta})$ com resíduos $r_i = y_i - \hat{y}_i(\vect{\theta})$:

\destaque{
\begin{equation*}
\vect{\theta}^{(k+1)} = \vect{\theta}^{(k)} + (\mathbf{J}^T\mathbf{J})^{-1}\mathbf{J}^T \vect{r}
\end{equation*}
}

\begin{columns}
\column{0.48\textwidth}
\textbf{Onde:}
\begin{itemize}\scriptsize
    \item $\mathbf{J}$: Jacobiana ($J_{ij} = \partial r_i / \partial \theta_j$)
    \item $\vect{r}$: vetor de resíduos individuais
\end{itemize}

\column{0.48\textwidth}
\begin{exampleblock}{Vantagens}\scriptsize
\begin{itemize}
    \item Evita o cálculo da Hessiana completa
    \item Excelente para mínimos quadrados
    \item Base do algoritmo Levenberg-Marquardt
\end{itemize}
\end{exampleblock}
\end{columns}
\end{frame}

\begin{frame}{Levenberg-Marquardt (LM): Algoritmo}
\definicao{Levenberg-Marquardt}{
Combina Gauss-Newton e gradient descent, interpolando entre ambos via parâmetro de amortecimento.
}

\destaque{
\begin{equation*}
\vect{\theta}^{(k+1)} = \vect{\theta}^{(k)} + (\mathbf{J}^T\mathbf{J} + \lambda \mathbf{I})^{-1}\mathbf{J}^T \vect{r}
\end{equation*}
}

\begin{itemize}\small
    \item $\lambda$: parâmetro de amortecimento (damping)
    \item $\lambda$ grande $\Rightarrow$ gradient descent (pequenos passos, robusto)
    \item $\lambda$ pequeno $\Rightarrow$ Gauss-Newton (passos grandes, rápido)
\end{itemize}
\end{frame}

\begin{frame}{Levenberg-Marquardt (LM): Estratégia}
\textbf{Estratégia Adaptativa de Ajuste:}
\begin{itemize}
    \item Se $J$ diminui: aceitar passo, diminuir $\lambda$ (mais agressivo)
    \item Se $J$ aumenta: rejeitar passo, aumentar $\lambda$ (mais conservador)
\end{itemize}

\vspace{0.5cm}

\importante{
LM é o método padrão para ajuste não-linear de curvas!
}
\end{frame}

\begin{frame}[fragile]{Otimização Local: Michaelis-Menten (curve\_fit)}
\begin{lstlisting}[language=Python, caption=Ajuste simples usando curve\_fit, basicstyle=\ttfamily\scriptsize]
from scipy.optimize import curve_fit
import numpy as np

# Modelo de Michaelis-Menten
def michaelis_menten(S, Vmax, Km):
    return Vmax * S / (Km + S)

# Dados experimentais
S_data = np.array([0.5, 1, 2, 5, 10, 20, 50])
v_data = np.array([0.4, 0.7, 1.2, 1.8, 2.1, 2.3, 2.4])

# Otimizacao via curve_fit
params, cov = curve_fit(michaelis_menten, S_data, v_data,
                         p0=[3, 5])
Vmax_fit, Km_fit = params

print(f"Vmax = {Vmax_fit:.3f}")
print(f"Km = {Km_fit:.3f}")
\end{lstlisting}
\end{frame}

\begin{frame}[fragile]{Otimização Local: Michaelis-Menten (least\_squares)}
\begin{lstlisting}[language=Python, caption=Ajuste via least\_squares (LM), basicstyle=\ttfamily\footnotesize]
from scipy.optimize import least_squares
import numpy as np

# Reuso dos dados anteriores e definicao de residuos
def residuals(params):
    Vmax, Km = params
    return v_data - michaelis_menten(S_data, Vmax, Km)

result = least_squares(residuals, x0=[3, 5], method='lm')
print(f"Parametros otimos: {result.x}")
print(f"Custo final: {result.cost:.6f}")
print(f"Sucesso: {result.success}")
\end{lstlisting}
\end{frame}

\begin{frame}[fragile]{Visualização do Ajuste}
\begin{lstlisting}[language=Python, caption=Comparação dados vs modelo ajustado, basicstyle=\ttfamily\scriptsize]
# Gerar curva do modelo ajustado
S_smooth = np.linspace(0, 60, 200)
v_fit = michaelis_menten(S_smooth, Vmax_fit, Km_fit)

# Plotar dados e ajuste
plt.figure(figsize=(10, 6))
plt.plot(S_data, v_data, 'ko', markersize=10, label='Dados exp')
plt.plot(S_smooth, v_fit, 'b-', linewidth=2,
         label=f'Ajuste: Vmax={Vmax_fit:.2f}, Km={Km_fit:.2f}')

# Linhas de referencia
plt.axhline(y=Vmax_fit, color='r', linestyle='--', alpha=0.5, label='Vmax')
plt.axvline(x=Km_fit, color='g', linestyle='--', alpha=0.5, label='Km')
plt.axhline(y=Vmax_fit/2, color='gray', linestyle=':')

plt.xlabel('[Substrato] (mM)'); plt.ylabel('Velocidade (umol/min)')
plt.title('Ajuste de Michaelis-Menten')
plt.legend(); plt.grid(True, alpha=0.3)
\end{lstlisting}
\end{frame}

% ============================================================================
% SEÇÃO 5: ALGORITMOS GLOBAIS
% ============================================================================
\section{Algoritmos Globais}

\begin{frame}{Por que Otimização Global?}
\begin{columns}
\column{0.48\textwidth}
\begin{alertblock}{Problema}
Métodos locais podem ficar presos em mínimos locais
\end{alertblock}

\begin{center}
\begin{tikzpicture}[scale=0.65]
    \draw[->] (0,0) -- (6.5,0) node[right] {$\theta$};
    \draw[->] (0,0) -- (0,3) node[above] {$J(\theta)$};

    \draw[thick, azulescuro, domain=0.4:6.0, samples=120, smooth]
        plot (\x, {1.2 + 0.8*sin(3.0*\x r) + 0.15*(\x-3.2)^2});

    \fill[red] (0.9,0.7) circle (2pt);
    \fill[red] (3.0,1.2) circle (2pt);
    \fill[verdeacido] (5.1,0.5) circle (3pt) node[below, font=\tiny] {Global};

    \node[above, font=\tiny] at (0.9,0.9) {Local 1};
    \node[above, font=\tiny] at (3.0,1.4) {Local 2};
\end{tikzpicture}
\end{center}

\column{0.48\textwidth}
\textbf{Quando Usar?}
\begin{itemize}\small
    \item Função objetivo altamente não-linear
    \item Múltiplos mínimos locais
    \item Sem boa estimativa inicial
    \item Parâmetros com grande incerteza
    \item Modelos biológicos complexos
\end{itemize}
\end{columns}
\end{frame}

\begin{frame}{Algoritmos Genéticos (GA)}
\definicao{Algoritmo Genético}{
Método inspirado na evolução biológica que mantém uma população de soluções candidatas e as evolui usando seleção, cruzamento e mutação.
}

\begin{columns}
\column{0.48\textwidth}
\textbf{Analogia Biológica:}
\begin{itemize}\small
    \item \textbf{Indivíduo:} parâmetros $\vect{\theta}$
    \item \textbf{Gene:} parâmetro individual $\theta_i$
    \item \textbf{Fitness:} inverso da função objetivo $1/J(\vect{\theta})$
    \item \textbf{População:} conjunto de soluções
\end{itemize}

\column{0.48\textwidth}
\begin{block}{Operadores Genéticos}\small
\begin{enumerate}
    \item \textbf{Seleção:} melhores soluções
    \item \textbf{Cruzamento:} combinar dois pais
    \item \textbf{Mutação:} variação aleatória
    \item \textbf{Elitismo:} manter o campeão
\end{enumerate}
\end{block}
\end{columns}
\end{frame}

\begin{frame}{Simulated Annealing (SA): Conceito}
\begin{block}{Inspiração}
Processo de recozimento em metalurgia: aquecer metal e esfriar lentamente para alcançar estrutura cristalina de mínima energia.
\end{block}

\textbf{Ideia Principal:}
\begin{itemize}
    \item Permitir movimentos para soluções piores com probabilidade decrescente.
    \item A "Temperatura" ($T$) controla a aceitação de soluções piores.
    \item Alta temperatura: exploração ampla do espaço.
    \item Baixa temperatura: refinamento e convergência local.
\end{itemize}
\end{frame}

\begin{frame}{Simulated Annealing (SA): Algoritmo}
\textbf{Probabilidade de Aceitação:}
\destaque{
\begin{equation*}
P(\text{aceitar}) = \begin{cases}
1 & \text{se } \Delta J < 0 \\
e^{-\Delta J / T} & \text{se } \Delta J \geq 0
\end{cases}
\end{equation*}
}
onde $\Delta J = J(\vect{\theta}_{\text{novo}}) - J(\vect{\theta}_{\text{atual}})$ e $T$ é a temperatura.

\vspace{0.5cm}

\begin{exampleblock}{Cronograma de Resfriamento}
$T_k = T_0 \cdot \alpha^k$ onde $0 < \alpha < 1$ (tipicamente $\alpha \in [0.8, 0.99]$)
\end{exampleblock}
\end{frame}

\begin{frame}{Differential Evolution (DE)}
\definicao{Evolução Diferencial}{
Algoritmo evolutivo que cria novos candidatos através de diferenças entre membros da população.
}

\begin{columns}
\column{0.48\textwidth}
\textbf{Operação Principal:}
\begin{enumerate}\small
    \item Escolher 3 indivíduos: $\vect{a}$, $\vect{b}$, $\vect{c}$
    \item Mutação: $\vect{v} = \vect{a} + F(\vect{b} - \vect{c})$
    \item Cruzamento usando taxa $CR$
    \item Seleção: manter o melhor
\end{enumerate}
\begin{itemize}\small
    \item Fator de escala $F \in [0.5, 1.0]$
    \item Crossover $CR \in [0.5, 0.9]$
\end{itemize}

\column{0.48\textwidth}
\begin{exampleblock}{Vantagens}\small
\begin{itemize}
    \item Simples de implementar
    \item Poucos parâmetros de ajuste
    \item Excelente robustez
    \item Não requer gradientes
    \item Ótimo para múltiplos mínimos
\end{itemize}
\end{exampleblock}
\end{columns}
\end{frame}

\begin{frame}{Particle Swarm Optimization (PSO)}
\begin{block}{Inspiração}
Comportamento social de bandos de pássaros ou cardumes de peixes.
\end{block}

\textbf{Conceito:}
\begin{itemize}\small
    \item Cada partícula tem posição $\vect{x}_i$ e velocidade $\vect{v}_i$.
    \item Movimento influenciado por: melhor posição própria ($\vect{p}_i$) e melhor posição global ($\vect{g}$).
\end{itemize}

\destaque{
\begin{align*}
\vect{v}_i^{(k+1)} &= w\vect{v}_i^{(k)} + c_1 r_1 (\vect{p}_i - \vect{x}_i^{(k)}) + c_2 r_2 (\vect{g} - \vect{x}_i^{(k)})\\
\vect{x}_i^{(k+1)} &= \vect{x}_i^{(k)} + \vect{v}_i^{(k+1)}
\end{align*}
}

\begin{itemize}\small
    \item $w$: peso de inércia | $c_1$, $c_2$: coeficientes cognitivo e social
    \item $r_1$, $r_2$: números aleatórios uniformes em $[0,1]$
\end{itemize}
\end{frame}

\begin{frame}[fragile]{Implementação: Differential Evolution (Função)}
\begin{lstlisting}[language=Python, caption=Modelo e funcao objetivo para DE, basicstyle=\ttfamily\footnotesize]
from scipy.optimize import differential_evolution
import numpy as np

# Funcao objetivo
def objective(params, t_data, y_data):
    """Modelo exponencial: y = a * exp(-b*t)"""
    a, b = params
    y_pred = a * np.exp(-b * t_data)
    sse = np.sum((y_data - y_pred)**2)
    return sse

# Dados experimentais
t_data = np.array([0, 1, 2, 3, 4, 5])
y_data = np.array([10, 6.1, 3.7, 2.2, 1.4, 0.8])
\end{lstlisting}
\end{frame}

\begin{frame}[fragile]{Implementação: Differential Evolution (Otimização)}
\begin{lstlisting}[language=Python, caption=Otimizacao global usando scipy, basicstyle=\ttfamily\scriptsize]
# Definir bounds para parametros: [(a_min, a_max), (b_min, b_max)]
bounds = [(1, 20), (0.01, 2.0)]

# Otimizar com DE
result = differential_evolution(
    objective, bounds, args=(t_data, y_data),
    strategy='best1bin', maxiter=1000, popsize=15,
    tol=1e-7, mutation=(0.5, 1), recombination=0.7, seed=42
)

print(f"Parametros otimos: a={result.x[0]:.3f}, b={result.x[1]:.3f}")
print(f"SSE minimo: {result.fun:.6f}")
print(f"Convergiu: {result.success}")
\end{lstlisting}
\end{frame}

\begin{frame}{Comparação de Métodos de Otimização}
\begin{center}
\begin{tabular}{lccccc}
\toprule
\textbf{Método} & \textbf{Velocidade} & \textbf{Global} & \textbf{Derivadas} & \textbf{Parâmetros} & \textbf{Uso} \\
\midrule
Grid Search & --- & Sim & Não & Poucos & Exploração \\
Gradient Descent & ++ & Não & Sim & Média & Local \\
Newton-Raphson & +++ & Não & 2ª ordem & Média & Refinamento \\
Levenberg-Marquardt & +++ & Não & Jacobiano & Média & Mín. Quad. \\
\midrule
GA & + & Sim & Não & Muitos & Robusto \\
Simulated Annealing & + & Sim & Não & Muitos & Explora \\
Differential Evolution & ++ & Sim & Não & Média & Melhor \\
Particle Swarm & ++ & Sim & Não & Média & Rápido \\
\bottomrule
\end{tabular}
\end{center}
\end{frame}

\begin{frame}{Diretrizes para Estimação de Parâmetros}
\begin{exampleblock}{Recomendações Práticas}
\begin{itemize}
    \item \textbf{Boa estimativa inicial + poucos parâmetros:} Use Levenberg-Marquardt (rápido e local).
    \item \textbf{Sem estimativa inicial + mínimos locais:} Use Differential Evolution (robusto e global).
    \item \textbf{Problema altamente não-linear:} Abordagem híbrida (rodar DE primeiro e usar resultado como inicial no LM).
    \item \textbf{Muitos parâmetros ($>20$):} Considere métodos Bayesianos (MCMC).
\end{itemize}
\end{exampleblock}
\end{frame}

% ============================================================================
% SEÇÃO 6: ANÁLISE DE IDENTIFICABILIDADE
% ============================================================================
\section{Análise de Identificabilidade}

\begin{frame}{Identificabilidade Estrutural: Exemplos}
\definicao{Identificabilidade Estrutural}{
Propriedade teórica do modelo: os parâmetros podem ser determinados unicamente a partir de dados perfeitos (infinitos e sem ruído)?
}

\begin{columns}
\column{0.48\textwidth}
\begin{block}{Identificável}
\begin{equation*}
\frac{dx}{dt} = k_1 - k_2 x
\end{equation*}
Ambos $k_1$ e $k_2$ podem ser determinados unicamente.
\end{block}

\column{0.48\textwidth}
\begin{alertblock}{Não Identificável}
\begin{equation*}
\frac{dx}{dt} = (k_1 k_2) - k_2 x
\end{equation*}
Apenas o produto $k_1 k_2$ é identificável, não $k_1$ e $k_2$ separadamente.
\end{alertblock}
\end{columns}
\end{frame}

\begin{frame}{Identificabilidade Estrutural: Diretrizes}
\importante{
Antes de estimar parâmetros, verifique se o modelo é estruturalmente identificável!
}

\vspace{0.5cm}

\textbf{Ferramentas de Análise:}
\begin{itemize}
    \item \textbf{Análise algébrica:} eliminação de variáveis via álgebra computacional.
    \item \textbf{Softwares Especializados:} DAISY, GenSSI, STRIKE-GOLDD.
\end{itemize}
\end{frame}

\begin{frame}{Identificabilidade Prática}
\definicao{Identificabilidade Prática}{
Podem os parâmetros ser determinados com precisão razoável a partir de dados experimentais reais (finitos e com ruído)?
}

\begin{columns}
\column{0.48\textwidth}
\textbf{Causas de Não-Identificabilidade:}
\begin{itemize}\small
    \item Dados insuficientes/baixa qualidade
    \item Parâmetros altamente correlacionados
    \item Modelo super-parametrizado (overfitting)
    \item Baixa sensibilidade a parâmetros
\end{itemize}

\column{0.48\textwidth}
\exemplo{Correlação de Parâmetros}{\small
Modelo: $y = a \cdot e^{-b \cdot t}$\\
Se $t$ é pequeno: $e^{-bt} \approx 1 - bt$\\
$\Rightarrow y \approx a(1 - bt) = a - abt$\\
Difícil separar $a$ de $ab$ (produto correlacionado).
}
\end{columns}
\end{frame}

\begin{frame}{Matriz de Informação de Fisher (FIM)}
\begin{block}{Quantifica Informação sobre Parâmetros}
Mede quanta informação os dados contêm sobre cada parâmetro.
\end{block}

\destaque{
\begin{equation*}
\mathbf{F}_{ij} = \sum_{k=1}^{n} \frac{1}{\sigma_k^2} \frac{\partial \hat{y}_k}{\partial \theta_i} \frac{\partial \hat{y}_k}{\partial \theta_j}
\end{equation*}
}

\begin{columns}
\column{0.48\textwidth}
\textbf{Onde:}
\begin{itemize}\small
    \item $\hat{y}_k$: previsão do modelo no ponto $k$
    \item $\sigma_k^2$: variância do erro no ponto $k$
    \item $\theta_i$, $\theta_j$: parâmetros
\end{itemize}

\column{0.48\textwidth}
\textbf{Interpretação:}
\begin{itemize}\small
    \item FIM grande $\Rightarrow$ alta precisão
    \item FIM pequena $\Rightarrow$ alta incerteza
    \item $\text{det}(\mathbf{F}) \approx 0$ $\Rightarrow$ não-identificável
\end{itemize}
\end{columns}
\end{frame}

\begin{frame}{Intervalos de Confiança: Teoria}
\begin{block}{Estimativa Assintótica}
Para grandes amostras, $\vect{\theta} \sim \mathcal{N}(\vect{\theta}^*, \mathbf{C})$ onde:
\end{block}

\destaque{
\begin{equation*}
\mathbf{C} = \text{Cov}(\vect{\theta}) \approx \mathbf{F}^{-1}
\end{equation*}
}

\textbf{Intervalo de Confiança (95\%):}
\begin{equation*}
\theta_i^* \pm 1.96 \sqrt{C_{ii}}
\end{equation*}
\end{frame}

\begin{frame}{Intervalos de Confiança: Aplicação Prática}
\begin{exampleblock}{Na Prática}
\begin{itemize}
    \item Diagonal: $\sqrt{C_{ii}}$ = erro padrão do parâmetro $i$
    \item Off-diagonal: $C_{ij}$ = correlação entre parâmetros $i$ e $j$
\end{itemize}
\end{exampleblock}

\importante{
Intervalos de confiança estreitos $\not\Rightarrow$ modelo correto!\\
Apenas indicam precisão do ajuste.
}
\end{frame}

\begin{frame}{Profile Likelihood}
\definicao{Profile Likelihood}{
Método robusto para calcular intervalos de confiança explorando o espaço de parâmetros sistematicamente.
}

\begin{columns}
\column{0.48\textwidth}
\textbf{Procedimento para $\theta_i$:}
\begin{enumerate}\small
    \item Fixar $\theta_i$ em valor específico
    \item Otimizar os outros parâmetros
    \item Registrar $J_{\min}(\theta_i)$
    \item Repetir para vários valores
    \item Plotar $J_{\min}$ vs $\theta_i$
\end{enumerate}

\column{0.48\textwidth}
\begin{center}
\begin{tikzpicture}[scale=0.7]
    \draw[->] (0,0) -- (5,0) node[right, font=\tiny] {$\theta_i$};
    \draw[->] (0,0) -- (0,2.5) node[above, font=\tiny] {$J_{\min}$};

    \draw[thick, azulescuro, domain=0.4:4.6, samples=80, smooth]
        plot (\x, {0.5*(\x-2.5)^2 + 0.3});

    \draw[dashed, red] (0,1.2) -- (5,1.2) node[right, font=\tiny] {Threshold};

    \draw[<->, thick, verdeacido] (1.2,0.1) -- (3.8,0.1);
    \node[below, font=\tiny] at (2.5,0.1) {IC 95\%};

    \fill[red] (2.5,0.3) circle (2pt) node[below, font=\tiny] {$\theta_i^*$};
\end{tikzpicture}
\end{center}
\end{columns}
\end{frame}

\begin{frame}[fragile]{Implementação: Visualização de Correlação}
\begin{lstlisting}[language=Python, caption=Heatmap de correlacao usando seaborn, basicstyle=\ttfamily\scriptsize]
from scipy.optimize import curve_fit
import seaborn as sns
import matplotlib.pyplot as plt

# Ajustar e extrair correlacao
params, cov = curve_fit(model_func, x_data, y_data)
std = np.sqrt(np.diag(cov))
corr = cov / np.outer(std, std)

# Plotar heatmap
param_names = ['k1', 'k2', 'k3']
plt.figure(figsize=(8, 6))
sns.heatmap(corr, annot=True, fmt='.2f',
            xticklabels=param_names, yticklabels=param_names,
            cmap='coolwarm', center=0, vmin=-1, vmax=1)
plt.title('Matriz de Correlacao')
\end{lstlisting}
\end{frame}

\begin{frame}[fragile]{Implementação: Alertas de Correlação}
\begin{lstlisting}[language=Python, caption=Deteccao de parametros colineares, basicstyle=\ttfamily\footnotesize]
# Alerta de alta correlacao
for i in range(len(params)):
    for j in range(i+1, len(params)):
        if abs(corr[i,j]) > 0.9:
            print(f"ALERTA: {param_names[i]} e {param_names[j]} "
                  f"altamente correlacionados (r={corr[i,j]:.2f})")
\end{lstlisting}
\end{frame}

\begin{frame}[fragile]{Implementação: Cálculo de Intervalos de Confiança}
\begin{lstlisting}[language=Python, caption=Calcular ICs usando distribuicao t, basicstyle=\ttfamily\footnotesize]
from scipy import stats
import numpy as np

params_best = params
errors = np.sqrt(np.diag(cov))

# Intervalo de confianca 95%
alpha = 0.05
t_val = stats.t.ppf(1-alpha/2, len(x_data)-len(params))
ci = t_val * errors
\end{lstlisting}
\end{frame}

\begin{frame}[fragile]{Implementação: Plot de Parâmetros e CIs}
\begin{lstlisting}[language=Python, caption=Visualizar parametros com barras de erro, basicstyle=\ttfamily\scriptsize]
import matplotlib.pyplot as plt

fig, ax = plt.subplots(figsize=(10, 6))
x_pos = np.arange(len(param_names))

ax.bar(x_pos, params_best, yerr=ci, alpha=0.7, capsize=10,
       color='steelblue', error_kw={'elinewidth': 2, 'capthick': 2})
ax.set_xticks(x_pos)
ax.set_xticklabels(param_names)
ax.set_ylabel('Valor do Parametro')
ax.set_title('Parametros Estimados com ICs 95%')
ax.grid(axis='y', alpha=0.3)

for name, val, err in zip(param_names, params_best, ci):
    print(f"{name} = {val:.4f} +/- {err:.4f} "
          f"({100*err/val:.1f}% erro relativo)")
\end{lstlisting}
\end{frame}

% ============================================================================
% SEÇÃO 7: VALIDAÇÃO E CRITÉRIOS DE AJUSTE
% ============================================================================
\section{Validação e Critérios de Ajuste}

\begin{frame}{Análise de Resíduos: Padrões}
\begin{block}{Verificação de Hipóteses}
Examinar resíduos $r_i = y_i - \hat{y}_i$ para validar suposições do modelo.
\end{block}

\textbf{Padrões a Verificar:}
\begin{columns}
\column{0.48\textwidth}
\begin{exampleblock}{Bom Ajuste}\small
\begin{itemize}
    \item Distribuição aleatória
    \item Sem padrão sistemático
    \item Centrados em zero
    \item Variância constante
    \item Aproximadamente normais
\end{itemize}
\end{exampleblock}

\column{0.48\textwidth}
\begin{alertblock}{Problemas}\small
\begin{itemize}
    \item Padrão curvo: inadequação
    \item Funil: heterocedasticidade
    \item Outliers: dados anômalos
    \item Não-normalidade do erro
\end{itemize}
\end{alertblock}
\end{columns}
\end{frame}

\begin{frame}{Análise de Resíduos: Visualização}
\begin{center}
\begin{tikzpicture}[scale=0.7]
    % Residual plot - good
    \begin{scope}
        \draw[->] (0,0) -- (4,0) node[right, font=\tiny] {$\hat{y}$};
        \draw[->] (0,-1) -- (0,1) node[above, font=\tiny] {$r$};
        \foreach \x in {0.5,1,1.5,2,2.5,3,3.5} {
            \pgfmathsetmacro{\y}{0.3*rand}
            \fill[azulescuro] (\x,\y) circle (1pt);
        }
        \node[below, font=\tiny] at (2,-1.3) {Bom (Aleatório)};
    \end{scope}

    % Residual plot - bad
    \begin{scope}[xshift=5cm]
        \draw[->] (0,0) -- (4,0) node[right, font=\tiny] {$\hat{y}$};
        \draw[->] (0,-1) -- (0,1) node[above, font=\tiny] {$r$};
        \draw[thick, red, domain=0.3:3.7, samples=50, smooth]
            plot (\x, {0.5*sin(3*\x r)});
        \foreach \x in {0.5,1,1.5,2,2.5,3,3.5} {
            \fill[red] (\x,{0.5*sin(3*\x r)}) circle (1pt);
        }
        \node[below, font=\tiny] at (2,-1.3) {Ruim (Padrão Curvo)};
    \end{scope}
\end{tikzpicture}
\end{center}
\end{frame}

\begin{frame}{Coeficiente de Determinação ($R^2$): Definição}
\definicao{$R^2$}{
Proporção da variância nos dados explicada pelo modelo. Varia de 0 (nenhum ajuste) a 1 (ajuste perfeito).
}

\destaque{
\begin{equation*}
R^2 = 1 - \frac{\text{SS}_{\text{res}}}{\text{SS}_{\text{tot}}} = 1 - \frac{\sum_i (y_i - \hat{y}_i)^2}{\sum_i (y_i - \bar{y})^2}
\end{equation*}
}

\begin{itemize}\small
    \item $\text{SS}_{\text{res}}$: soma dos quadrados dos resíduos
    \item $\text{SS}_{\text{tot}}$: soma dos quadrados totais
    \item $\bar{y}$: média dos dados experimentais
\end{itemize}
\end{frame}

\begin{frame}{Coeficiente de Determinação ($R^2$): Limitações}
\begin{alertblock}{Limitações Importantes}
\begin{itemize}
    \item $R^2$ sempre aumenta ao adicionar novos parâmetros (mesmo irrelevantes).
    \item Um $R^2$ alto não garante que o modelo é correto ou preditivo.
    \item Pode ser altamente enganoso para modelos não-lineares.
    \item Use o $R^2$ ajustado para comparar modelos com diferentes números de parâmetros.
\end{itemize}
\end{alertblock}
\end{frame}

\begin{frame}{$R^2$ Ajustado: Definição}
\begin{block}{Penaliza Complexidade do Modelo}
Ajusta o $R^2$ baseando-se no número de graus de liberdade.
\end{block}

\destaque{
\begin{equation*}
R^2_{\text{adj}} = 1 - \frac{(1-R^2)(n-1)}{n-p-1}
\end{equation*}
}

onde:
\begin{itemize}\small
    \item $n$: número de observações (tamanho da amostra)
    \item $p$: número de parâmetros estimáveis
\end{itemize}

\vspace{0.3cm}

\textbf{Características:}
\begin{itemize}\small
    \item Pode diminuir ao adicionar parâmetros ruins ou irrelevantes.
    \item Mais apropriado para comparação de modelos.
\end{itemize}
\end{frame}

\begin{frame}{Akaike Information Criterion (AIC): Teoria}
\definicao{AIC}{
Critério que balanceia a qualidade do ajuste com a complexidade do modelo.
}

\destaque{\small
\begin{equation*}
\text{AIC} = 2p - 2\ln(\mathcal{L}) \approx n\ln\left(\frac{\text{SSE}}{n}\right) + 2p
\end{equation*}
}

\begin{columns}
\column{0.48\textwidth}
\textbf{Onde:}
\begin{itemize}\scriptsize
    \item $p$: número de parâmetros
    \item $\mathcal{L}$: máxima verossimilhança
    \item SSE: soma dos resíduos quadrados
\end{itemize}

\vspace{0.1cm}
\textbf{Interpretação:}
\begin{itemize}\scriptsize
    \item Menor AIC = melhor modelo
    \item Penaliza complexidade extra
    \item $\Delta\text{AIC} > 10$: evidência forte
\end{itemize}

\column{0.48\textwidth}
\begin{exampleblock}{AIC Corrigido (AICc)}\small
Para amostras pequenas ($n/p < 40$):
\begin{equation*}
\text{AICc} = \text{AIC} + \frac{2p(p+1)}{n-p-1}
\end{equation*}
\end{exampleblock}
\end{columns}
\end{frame}

\begin{frame}{Bayesian Information Criterion (BIC): Definição}
\definicao{BIC (ou Schwarz Criterion)}{
Similar ao AIC, mas penaliza mais fortemente a complexidade do modelo.
}

\destaque{
\begin{equation*}
\text{BIC} = p\ln(n) - 2\ln(\mathcal{L}) \approx n\ln\left(\frac{\text{SSE}}{n}\right) + p\ln(n)
\end{equation*}
}

\textbf{Comparação AIC vs BIC:}
\begin{itemize}\small
    \item BIC penaliza complexidade mais forte (termo $p\ln(n)$ vs $2p$ do AIC).
    \item Para $n > 7$, o BIC seleciona modelos mais parcimoniosos (menores).
    \item O AIC busca otimizar a previsão futura, enquanto o BIC tenta identificar o "modelo verdadeiro".
\end{itemize}
\end{frame}

\begin{frame}{Comparação de Critérios de Seleção (AIC vs BIC)}
\begin{center}
\begin{tabular}{lcc}
\toprule
\textbf{Critério} & \textbf{Objetivo} & \textbf{Preferência} \\
\midrule
AIC & Previsão / Incerteza & Modelos maiores \\
BIC & Identificação do Processo & Modelos menores \\
\bottomrule
\end{tabular}
\end{center}

\vspace{0.5cm}

\begin{exampleblock}{Resumo Prático}
Na biologia de sistemas, onde os dados costumam ser limitados e o ruído é alto, o **AICc** e o **BIC** são usados em conjunto. Se divergirem, prioriza-se o modelo mais simples (BIC) para evitar overfitting.
\end{exampleblock}
\end{frame}

\begin{frame}{Cross-Validation (Validação Cruzada)}
\begin{block}{Teste de Capacidade Preditiva}
Avaliar o desempenho do modelo em dados não vistos durante o treinamento para evitar overfitting.
\end{block}

\begin{columns}
\column{0.48\textwidth}
\textbf{K-Fold Cross-Validation:}
\begin{enumerate}\small
    \item Dividir dados em $K$ folds.
    \item Para cada fold $k$:
    \begin{itemize}
        \item Treinar modelo nos outros $K-1$ folds.
        \item Testar no fold $k$.
        \item Registrar erro de previsão.
    \end{itemize}
    \item Calcular erro médio de CV.
\end{enumerate}

\column{0.48\textwidth}
\begin{center}
\begin{tikzpicture}[scale=0.65]
    \foreach \i in {0,...,4} {
        \pgfmathsetmacro{\x}{\i*1.1}
        \ifnum\i=2
            \draw[fill=red!40] (\x,0) rectangle +(0.9,0.4);
            \node at (\x+0.45,0.2) {\tiny Test};
        \else
            \draw[fill=azulclaro!40] (\x,0) rectangle +(0.9,0.4);
            \node at (\x+0.45,0.2) {\tiny Train};
        \fi
    }
    \node[below, font=\tiny] at (2.5,-0.2) {5-Fold CV (Dobra 3 como teste)};
\end{tikzpicture}
\end{center}

\begin{block}{Diretrizes}\small
Típico: $K=5$ ou $K=10$. Se $K=n$, chamamos de Leave-One-Out (LOO), muito custoso.
\end{block}
\end{columns}
\end{frame}

\begin{frame}{Overfitting vs Underfitting}
\begin{center}
\begin{tikzpicture}[scale=0.75]
    % Underfitting
    \begin{scope}
        \draw[->] (0,0) -- (3,0) node[right, font=\tiny] {$x$};
        \draw[->] (0,0) -- (0,2.2) node[above, font=\tiny] {$y$};
        \foreach \x/\y in {0.5/0.8, 1/1.3, 1.5/1.7, 2/1.4, 2.5/0.9} {
            \fill[blue] (\x,\y) circle (1.5pt);
        }
        \draw[thick, red] (0.3,1.3) -- (2.7,1.3);
        \node[below, font=\tiny] at (1.5,-0.2) {Underfitting};
        \node[below, font=\tiny] at (1.5,-0.5) {Viés alto};
    \end{scope}

    % Good fit
    \begin{scope}[xshift=3.8cm]
        \draw[->] (0,0) -- (3,0) node[right, font=\tiny] {$x$};
        \draw[->] (0,0) -- (0,2.2) node[above, font=\tiny] {$y$};
        \foreach \x/\y in {0.5/0.8, 1/1.3, 1.5/1.7, 2/1.4, 2.5/0.9} {
            \fill[blue] (\x,\y) circle (1.5pt);
        }
        \draw[thick, verdeacido, domain=0.3:2.7, samples=50, smooth]
            plot (\x, {0.8 + 1.8*sin((\x-0.5)*120)});
        \node[below, font=\tiny] at (1.5,-0.2) {Bom Ajuste};
        \node[below, font=\tiny] at (1.5,-0.5) {Balanceado};
    \end{scope}

    % Overfitting
    \begin{scope}[xshift=7.6cm]
        \draw[->] (0,0) -- (3,0) node[right, font=\tiny] {$x$};
        \draw[->] (0,0) -- (0,2.2) node[above, font=\tiny] {$y$};
        \foreach \x/\y in {0.5/0.8, 1/1.3, 1.5/1.7, 2/1.4, 2.5/0.9} {
            \fill[blue] (\x,\y) circle (1.5pt);
        }
        \draw[thick, red, samples=100, smooth]
            plot coordinates {(0.5,0.8) (1,1.3) (1.5,1.7) (2,1.4) (2.5,0.9)};
        \node[below, font=\tiny] at (1.5,-0.2) {Overfitting};
        \node[below, font=\tiny] at (1.5,-0.5) {Variância alta};
    \end{scope}
\end{tikzpicture}
\end{center}

\begin{alertblock}{Sintomas de Overfitting}
\begin{itemize}\small
    \item Erro de treinamento muito menor que o erro de teste.
    \item Modelo complexo com excesso de parâmetros livres.
    \item ICs excessivamente largos; previsões ruins fora do domínio.
\end{itemize}
\end{alertblock}
\end{frame}

\begin{frame}[fragile]{Implementação: Cálculo de Métricas (RMSE, R2)}
\begin{lstlisting}[language=Python, caption=Calculo de RMSE e R-squared ajustado, basicstyle=\ttfamily\scriptsize]
import numpy as np

def calculate_metrics(y_true, y_pred, n_params):
    n = len(y_true)
    residuals = y_true - y_pred

    # SSE e MSE
    sse = np.sum(residuals**2)
    mse = sse / n
    rmse = np.sqrt(mse)

    # R-squared e R-squared ajustado
    ss_tot = np.sum((y_true - np.mean(y_true))**2)
    r2 = 1 - (sse / ss_tot)
    r2_adj = 1 - (1-r2)*(n-1)/(n-n_params-1)
\end{lstlisting}
\end{frame}

\begin{frame}[fragile]{Implementação: Cálculo de Métricas (AIC, BIC)}
\begin{lstlisting}[language=Python, caption=Calculo de AIC/BIC e execucao, basicstyle=\ttfamily\scriptsize]
    # Calculo continuacao: AIC e BIC
    aic = n*np.log(sse/n) + 2*n_params
    bic = n*np.log(sse/n) + n_params*np.log(n)

    return {
        'SSE': sse, 'MSE': mse, 'RMSE': rmse,
        'R2': r2, 'R2_adj': r2_adj, 'AIC': aic, 'BIC': bic,
        'residuals': residuals
    }

# Executar e exibir resultados
metrics = calculate_metrics(y_data, y_pred, n_params=2)
for key, val in metrics.items():
    if key != 'residuals':
        print(f"{key}: {val:.4f}")
\end{lstlisting}
\end{frame}

% ============================================================================
% SEÇÃO 8: EXERCÍCIOS
% ============================================================================
\section{Exercícios}

\begin{frame}{Exercícios - Parte 1: Busca em Grade}
\begin{block}{Questão 1: Busca em Grade}
Implemente busca em grade para encontrar parâmetros $a$ e $b$ do modelo $y = a \cdot x^b$ usando os dados:
\begin{center}
\begin{tabular}{c|cccccc}
$x$ & 1 & 2 & 3 & 4 & 5 & 6 \\
\hline
$y$ & 2.1 & 8.3 & 18.2 & 32.5 & 50.1 & 72.3
\end{tabular}
\end{center}
Use: $a \in [1, 3]$, $b \in [1.5, 2.5]$, grade 20x20.
\end{block}
\end{frame}

\begin{frame}{Exercícios - Parte 1: Descida do Gradiente}
\begin{block}{Questão 2: Gradient Descent}
Implemente gradient descent manualmente para minimizar $J(\theta) = (\theta - 5)^2 + 3$:
\begin{enumerate}
    \item Use $\theta_0 = 0$ e $\alpha = 0.1$
    \item Plote trajetória de convergência
    \item Teste diferentes valores de $\alpha$ (0.01, 0.1, 0.5, 1.0)
    \item Compare número de iterações até convergência
\end{enumerate}
\end{block}
\end{frame}

\begin{frame}{Exercícios - Parte 2: Levenberg-Marquardt}
\begin{block}{Questão 3: Levenberg-Marquardt}
Ajuste parâmetros da cinética de Michaelis-Menten $v = \frac{V_{\max}[S]}{K_M + [S]}$ aos dados:
\begin{center}
\begin{tabular}{c|ccccccc}
$[S]$ (mM) & 0.1 & 0.5 & 1 & 2 & 5 & 10 & 20 \\
\hline
$v$ ($\mu$mol/min) & 0.3 & 1.1 & 1.8 & 2.5 & 3.3 & 3.7 & 3.9
\end{tabular}
\end{center}
\begin{enumerate}
    \item Use \texttt{scipy.optimize.curve\_fit}
    \item Calcule intervalos de confiança
    \item Plote dados vs modelo ajustado e faça análise de resíduos
\end{enumerate}
\end{block}
\end{frame}

\begin{frame}{Exercícios - Parte 2: Differential Evolution}
\begin{block}{Questão 4: Differential Evolution}
Compare DE com LM no problema anterior:
\begin{enumerate}
    \item Use bounds amplos: $V_{\max} \in [0.1, 10]$, $K_M \in [0.1, 10]$
    \item Qual método é mais rápido? Qual encontra melhor solução?
    \item Use DE para encontrar boa estimativa inicial para o LM
\end{enumerate}
\end{block}
\end{frame}

\begin{frame}{Exercícios - Parte 3: Modelo Logístico}
\begin{alertblock}{Questão 5: Crescimento Logístico}
Considere dados de crescimento bacteriano:
\begin{enumerate}\small
    \item Gere dados sintéticos: $N(t) = \frac{K}{1 + \left(\frac{K}{N_0}-1\right)e^{-rt}}$ com $N_0=1, r=0.5, K=100, t \in [0, 20]$ + ruído.
    \item Estime $r$ e $K$ usando três métodos diferentes.
    \item Analise identificabilidade e profile likelihood.
\end{enumerate}
\end{alertblock}
\end{frame}

\begin{frame}{Exercícios - Parte 3: Validação Cruzada}
\begin{block}{Questão 6: Validação Cruzada}
Implemente 5-fold cross-validation para o modelo logístico:
\begin{enumerate}
    \item Divida dados em 5 folds.
    \item Calcule RMSE médio de CV.
    \item Compare com RMSE de treinamento (checar overfitting).
    \item Teste com modelos de complexidade crescente.
\end{enumerate}
\end{block}
\end{frame}

\begin{frame}{Exercícios - Parte 4: Descrição do Projeto}
\begin{alertblock}{Projeto: Sistema de Lotka-Volterra}
Estime parâmetros do sistema predador-presa:
\begin{align*}
\frac{dV}{dt} &= \alpha V - \beta V P\\
\frac{dP}{dt} &= \delta \beta V P - \gamma P
\end{align*}
\textbf{Dados simulados:} Gere séries temporais com parâmetros conhecidos + ruído gaussiano.
\end{alertblock}
\end{frame}

\begin{frame}{Exercícios - Parte 4: Tarefas do Projeto}
\begin{alertblock}{Tarefas Recomendadas}
\begin{enumerate}
    \item Implementar função de custo (SSE de ambas variáveis).
    \item Estimar 4 parâmetros usando Differential Evolution.
    \item Refinar estimativas com Levenberg-Marquardt.
    \item Analisar identificabilidade (matriz de correlação de Fisher).
    \item Calcular intervalos de confiança assintóticos e plotá-los.
    \item Validar o ajuste simulando o modelo vs dados ruidosos.
\end{enumerate}
\end{alertblock}
\end{frame}

% ============================================================================
% SEÇÃO 9: REFERÊNCIAS
% ============================================================================
\section{Referências}

\begin{frame}{Referências Principais}
\begin{thebibliography}{99}
\bibitem{nocedal} Nocedal J., Wright S.J. (2006). \textit{Numerical Optimization}, 2nd ed. Springer.

\bibitem{raue} Raue A. et al. (2009). Structural and practical identifiability analysis of partially observed dynamical models. \textit{Bioinformatics}.

\bibitem{voit} Voit E.O. et al. (2015). 150 years of the mass action law. \textit{PLoS Computational Biology}.

\bibitem{moles} Moles C.G. et al. (2003). Parameter estimation in biochemical pathways. \textit{Genome Research}.

\bibitem{burnham} Burnham K.P., Anderson D.R. (2002). \textit{Model Selection and Multimodel Inference}, 2nd ed. Springer.
\end{thebibliography}
\end{frame}

\begin{frame}{Leitura Complementar}
\begin{columns}
\column{0.48\textwidth}
\begin{block}{Livros}\scriptsize
\begin{itemize}
    \item Press W.H. et al. (2007). \textit{Numerical Recipes}. Cambridge.
    \item Bard Y. (1974). \textit{Nonlinear Parameter Estimation}.
    \item Seber G.A.F., Wild C.J. (2003). \textit{Nonlinear Regression}. Wiley.
\end{itemize}
\end{block}

\column{0.48\textwidth}
\begin{block}{Artigos Importantes}\scriptsize
\begin{itemize}
    \item Gutenkunst R.N. (2007). Universally sloppy parameter sensitivities.
    \item Villaverde A.F. (2016). Cooperative parameter estimation.
\end{itemize}
\end{block}

\begin{block}{Software e Tutoriais}\scriptsize
\begin{itemize}
    \item \textbf{SciPy:} \url{https://docs.scipy.org/doc/scipy/reference/optimize.html}
    \item \textbf{lmfit:} \url{https://lmfit.github.io}
    \item \textbf{PyMC:} \url{https://www.pymc.io}
\end{itemize}
\end{block}
\end{columns}
\end{frame}

% ============================================================================
% SLIDE FINAL
% ============================================================================
\begin{frame}
\begin{center}
{\Huge Obrigado!}

\vspace{1cm}

{\Large Capítulo 5: Estimação de Parâmetros}

\vspace{1cm}

\textbf{Próximo Capítulo:}\\
Sistemas Gênicos

\vspace{1cm}

\textit{Dúvidas? Perguntas?}
\end{center}
\end{frame}

\end{document}
